\documentclass[12pt]{article}
\usepackage[margin=1in]{geometry}
\usepackage{times}           % Times New Roman body text
\usepackage{graphicx,booktabs,amsmath,amssymb,amsthm,hyperref}
\usepackage[font=small]{caption}
\newtheorem{proposition}{Proposition}
\newtheorem{remark}{Remark}
\title{A Virtual Yeast Cell for Metabolic Engineering:\\
Genome-Scale FBA plus Graph and Sequence Deep Learning\\
for Gene Essentiality and Strain Design}
\author{MEGA-PROGRAM-27, Item 5}
\date{\today}
\begin{document}
\maketitle
\begin{abstract}
We combine the consensus yeast genome-scale metabolic model (yeast-GEM,
4{,}105 reactions / 2{,}748 metabolites / 1{,}143 genes) with a graph
convolutional network over the metabolic gene graph and a convolutional
sequence model to predict gene essentiality, benchmarked against the
published FBA knockout rule under a locked, gene-held-out cross-validation
protocol, and we apply the same model to strain design for succinate
overproduction, validated against published experimental deletions.
\end{abstract}

\section{Introduction}
% genome-scale models; FBA essentiality benchmark (acc 0.9015, MCC 0.532,
% 94/159 essentials missed); hypothesis: network + sequence context recovers
% FBA-missed essentials without reading labels.

\section{Methods}
\subsection{Model and media}
% yeast-GEM SBML in-repo, sha256 recorded; Kennedy Y7 medium reproducing the
% published benchmark; solver and reproducibility.
\subsection{Flux balance analysis}
\begin{align}
\max_{v}\quad & c^{\top} v \\
\text{s.t.}\quad & S v = 0, \quad l \le v \le u .
\end{align}
\subsection{Label-free mechanistic features}
% GPR isozyme redundancy, pFBA flux carriage, KO growth ratio, topology.
\subsection{Graph convolutional network}
\begin{equation}
H^{(l+1)} = \mathrm{ReLU}\!\left(\hat{D}^{-1/2}\hat{A}\hat{D}^{-1/2} H^{(l)} W^{(l)}\right)
\end{equation}
\subsection{Sequence CNN}
\subsection{Locked evaluation protocol}
% repeated stratified 5-fold gene-held-out CV; shared folds; train-only
% scaling and thresholds; paired bootstrap CI on delta-MCC.
\subsection{Strain design}
% succinate production envelope; single-KO scan at a growth floor;
% validation set from Raab et al. 2010 (SDH1, SDH2, IDH1, IDP1).

\section{Mathematical results}
\begin{proposition}[Weak duality bound for growth]
% shadow prices; dual of FBA LP; proof.
\end{proposition}
\begin{proposition}[Envelope monotonicity]
% max product flux is nonincreasing in the enforced growth fraction; proof
% by feasible-set nesting.
\end{proposition}

\section{Results}
\subsection{FBA benchmark reproduction}
\subsection{Cross-validated essentiality prediction}
\subsection{Strain-design validation}

\section{Discussion}
\section{Conclusion}
\bibliographystyle{plain}
\begin{thebibliography}{9}
\bibitem{raab2010} Raab et al., Metab.\ Eng.\ 12(6):518--525 (2010).
\bibitem{kipf2017} Kipf \& Welling, ICLR (2017).
\bibitem{yeastgem} Lu et al., Nat.\ Commun.\ 10:2390 (2019).
\end{thebibliography}
\end{document}
